\documentclass{article}

\usepackage{amsmath,amssymb}
\usepackage{xurl}
\usepackage[hidelinks]{hyperref}

\input{command}

\title{Wolf Proposition~1.5: Factor and Coordinate Direct-Sum Scope Restrictions}
\author{}
\date{2026-07-18}

\begin{document}
\maketitle
\gapnote{scope-restriction}{resolved}

\section{Source statement}

Wolf's Proposition~1.5 classifies positive linear retractions onto arbitrary
unital finite-dimensional $*$-subalgebras
\cite[Proposition~1.5]{Wolf2012Quantum}. After unitary conjugation, such a
subalgebra has the block form
\begin{align}
    \mathcal A
        &=\bigoplus_{k=1}^{K}
          M_{d_k}(\mathbb C)\otimes\one_{m_k}.
    \label{eq:wolf_prop1_5_one_factor_scope_source_block_algebra}
\end{align}
The proposition states that a positive retraction onto $\mathcal A$ is a
direct sum of weighted partial traces, one on each summand; see source
Equation~(1.40) \cite[Equation~(1.40)]{Wolf2012Quantum}. Its proof first shows
that the retraction removes inputs between distinct central summands and
produces no output in an incorrect summand. It then proves the
weighted-partial-trace formula on each factor separately. The local source is
\path{Notes/WolfNoteTexSource/ch01_deconstructing_quantum.tex},
lines~536--570.

\section{The one-factor theorem}

The declaration
\leanid{Matrix.exists_density_of_positive_retraction_onto_right_factor}
proves the complete factor argument for a positive retraction whose range is
\begin{align}
    \one_m\otimes M_d(\mathbb C).
    \label{eq:wolf_prop1_5_one_factor_scope_formal_factor_range}
\end{align}
It introduces no hypothesis absent from the one-factor specialization of
Proposition~1.5. In particular, the representing matrix $\rho$ is required
only to satisfy
\begin{align}
    \rho&\geq0,
    \label{eq:wolf_prop1_5_one_factor_scope_density_positive}\\
    \tr(\rho)&=1.
    \label{eq:wolf_prop1_5_one_factor_scope_density_trace_one}
\end{align}
The theorem therefore gives a faithful one-factor case of the source result,
but by itself it does not establish the direct-sum classification in
Eq.~\ref{eq:wolf_prop1_5_one_factor_scope_source_block_algebra}.

\section{The coordinate direct-sum construction}

The declaration
\leanid{Matrix.coordinateDirectSumConditionalExpectation_isKrausCPTP_and_isConditionalExpectation}
constructs the normalized trace-preserving choice simultaneously on all
summands in displayed direct-sum coordinates. In QICLean's tensor-factor
order, its range is
\begin{align}
    \bigoplus_{k=1}^{K}
        \one_{m_k}\otimes M_{d_k}(\mathbb C).
    \label{eq:wolf_prop1_5_one_factor_scope_coordinate_range}
\end{align}
This order is obtained from
Eq.~\ref{eq:wolf_prop1_5_one_factor_scope_source_block_algebra} by
interchanging the two tensor factors in each summand.

The coordinate theorem performs the coordinatewise partial traces and removes
the off-diagonal blocks. It does not, by itself, conjugate an arbitrarily
represented subalgebra into the coordinates of
Eq.~\ref{eq:wolf_prop1_5_one_factor_scope_coordinate_range}, nor does it
corestrict the ambient matrix-valued map to that subalgebra. It is therefore a
coordinate construction rather than the full unitarily transported
classification of Proposition~1.5 \cite{Wolf2012Quantum}.

\section{Resolution of the direct-sum argument}

The declaration
\leanid{StarSubalgebra.exists_block_densities_of_positive_retraction} combines
the one-factor theorem with the finite-dimensional block decomposition. Its
proof proceeds as follows:
\begin{enumerate}
    \item It conjugates and reindexes the ambient matrices using the unitary
          block data.
    \item It proves that the positive retraction annihilates inputs between
          distinct central summands.
    \item It proves that an input supported in one summand has no output in a
          different summand.
    \item It applies
          \leanid{Matrix.exists_density_of_positive_retraction_onto_right_factor}
          to each diagonal summand.
    \item It conjugates the resulting direct sum of weighted partial traces
          back to the original subalgebra.
\end{enumerate}

The second and third steps formalize the off-diagonal positivity arguments
that precede source Equation~(1.42)
\cite[Equation~(1.42)]{Wolf2012Quantum}. For the central projection $p_k$ of
the $k$th summand, these arguments use the complementary projection
\begin{align}
    \one-p_k.
    \label{eq:wolf_prop1_5_one_factor_scope_complementary_central_projection}
\end{align}
They eliminate cross-summand inputs and wrong-summand outputs before the
factor theorem is applied.

The resulting theorem supplies positive trace-one block densities $\rho_k$
and reconstructs the source classification as a unitarily transported direct
sum of weighted partial traces. The distinction between classifying a given
retraction and constructing a fixed-point projection later in the development
remains recorded in
\path{docs/paper-gaps/wolf_theorem6_14_fixed_point_projection_gap.tex}.

\section{Classification}

The one-factor theorem and the coordinate direct-sum construction were
individually scope-restricted relative to Wolf's Proposition~1.5. The theorem
\leanid{StarSubalgebra.exists_block_densities_of_positive_retraction} now
combines the missing block, support, and unitary-transport arguments. The
scope restriction is therefore resolved, and the formal development covers
the full finite-dimensional classification stated in Proposition~1.5
\cite[Proposition~1.5]{Wolf2012Quantum}.

\bibliographystyle{alpha}
\bibliography{references}

\end{document}
